\section*{Capítulo \ref{ch:g12:proton-nmr} --- RMN del protón}

\begin{solution}{exo:g12:proton-nmr:1}
Metano: 1. Etano: 1 (los dos \ce{CH3} son iguales). Propano: 2 (dos
\ce{CH3} iguales y un \ce{CH2}). Metanol: 2 (\ce{CH3} y \ce{OH}).
Metoximetano: 1.
\end{solution}

\begin{solution}{exo:g12:proton-nmr:2}
2 vecinos: un triplete, 1:2:1. 3 vecinos: un cuadruplete, 1:3:3:1.
0 vecinos: un singlete.
\end{solution}

\begin{solution}{exo:g12:proton-nmr:3}
Total \qty{40}{mm} para 8 \omterm{def:g9:inside-the-atom:proton}{protones}: \qty{5}{mm} por \omterm{def:g9:inside-the-atom:proton}{protón}. Las señales
representan 3, 2 y 3 \omterm{def:g9:inside-the-atom:proton}{protones}.
\end{solution}

\begin{solution}{exo:g12:proton-nmr:4}
\omterm{def:g11:functional-groups:carbonyl}{Aldehído}: de 9.7 a \qty{10.0}{ppm}. Un \ce{CH3} de una cadena: de 0.7 a
\qty{1.3}{ppm}.
\end{solution}

\begin{solution}{exo:g12:proton-nmr:5}
Sus doce \omterm{def:g12:proton-nmr:equivalent}{protones equivalentes} dan una sola línea, fina e intensa, fácil
de situar en el cero. El \omterm{def:g6:solutions-and-solubility:solute}{disolvente} deuterado no da ninguna señal de
\omterm{def:g9:inside-the-atom:proton}{protón} propia, que de lo contrario taparía las de los pocos miligramos
de muestra.
\end{solution}

\begin{solution}{exo:g12:proton-nmr:6}
Dos señales. \ce{CH2} vecino del cloro: 2 \omterm{def:g9:inside-the-atom:proton}{protones}, un cuadruplete (3
vecinos), en el intervalo 2.5--\qty{4.0}{ppm}. \ce{CH3}: 3 \omterm{def:g9:inside-the-atom:proton}{protones}, un
triplete (2 vecinos), cerca del intervalo habitual de los \ce{CH3}, algo
más arriba porque el cloro está a dos \omterm{def:g7:atoms-and-molecules:atom}{átomos} de distancia.
\end{solution}

\begin{solution}{exo:g12:proton-nmr:7}
Tres señales: los dos \ce{CH3}, 6 \omterm{def:g12:proton-nmr:equivalent}{protones equivalentes}, un doblete (1
vecino, el \ce{CH}); el \ce{CH}, 1 \omterm{def:g9:inside-the-atom:proton}{protón}, desdoblado por 6 vecinos en
7 líneas, en el intervalo \ce{H-C-O} (3.3--\qty{4.5}{ppm}); el
\ce{OH}, 1 \omterm{def:g9:inside-the-atom:proton}{protón}, un singlete.
\end{solution}

\begin{solution}{exo:g12:proton-nmr:8}
(a) propanona; (b) etanoato de etilo; (c) ácido propanoico.
\end{solution}

\begin{solution}{exo:g12:proton-nmr:9}
El \ce{CH2} está unido al \omterm{def:g7:atoms-and-molecules:atom}{átomo} de oxígeno del \omterm{def:g11:functional-groups:carboxylic-acid}{éster}, que es
electronegativo y desapantalla sus \omterm{def:g9:inside-the-atom:proton}{protones}: intervalo \ce{H-C-O},
3.3--\qty{4.5}{ppm}. El \ce{CH3} está un carbono más lejos y se queda en
el intervalo de las cadenas.
\end{solution}

\begin{solution}{exo:g12:proton-nmr:10}
La propanona: un \ce{C=O} (\omterm{def:g11:functional-groups:carbonyl}{cetona}, \qty{1715}{cm^{-1}}) y seis \omterm{def:g12:proton-nmr:equivalent}{protones
equivalentes}, un solo singlete. El propanal daría tres señales, una de
ellas, la del \omterm{def:g9:inside-the-atom:proton}{protón} aldehídico de \ce{CH3-CH2-CHO}, cerca de 9.7--\qty{10.0}{ppm}.
\end{solution}

\begin{solution}{exo:g12:proton-nmr:11}
El \omterm{def:g9:inside-the-atom:proton}{protón} del \ce{OH} se intercambia rápidamente entre \omterm{def:g7:atoms-and-molecules:molecule}{moléculas}: da un
singlete y no desdobla a sus vecinos. Por tanto, el \ce{CH2} solo está
desdoblado por los 3 \omterm{def:g9:inside-the-atom:proton}{protones} del \ce{CH3}: un cuadruplete.
\end{solution}

\begin{solution}{exo:g12:proton-nmr:12}
Ácido butanoico \ce{CH3-CH2-CH2-COOH}: 4 señales, un triplete (\ce{CH3}),
una señal de 6 líneas (\ce{CH2} central, 5 vecinos), un triplete
(\ce{CH2-C=O}) y un singlete cerca de 11--\qty{12}{ppm}. Etanoato de
etilo: un singlete, un cuadruplete y un triplete. Propanoato de metilo
\ce{CH3-CH2-CO-O-CH3}: un triplete, un cuadruplete (\ce{CH2-C=O},
2.0--\qty{2.4}{ppm}) y un singlete (\ce{O-CH3}). Metanoato de propilo
\ce{H-CO-O-CH2-CH2-CH3}: 4 señales, un singlete muy a la izquierda (el H
del carbono del \ce{C=O}), un triplete (\ce{O-CH2}), una señal de 6
líneas y un triplete (\ce{CH3}).
\end{solution}

\begin{solution}{exo:g12:proton-nmr:13}
La señal del \ce{OH}: los \omterm{def:g9:inside-the-atom:proton}{protones} \ce{O-H} se intercambian con el
deuterio del agua pesada, y el \ce{O-D} no da señal de \omterm{def:g9:inside-the-atom:proton}{protón}. Una señal
que desaparece al agitar con \ce{D2O} corresponde a un \omterm{def:g9:inside-the-atom:proton}{protón} unido a O
o a N.
\end{solution}

\begin{solution}{exo:g12:proton-nmr:14}
El etanol añade un cuadruplete cerca de \qty{3.69}{ppm}, un triplete
cerca de \qty{1.23}{ppm} y un singlete del \ce{OH}. El singlete del
\omterm{def:g11:functional-groups:carboxylic-acid}{éster}: \qty{30}{mm} para 3 \omterm{def:g9:inside-the-atom:proton}{protones}, \qty{10}{mm} por \omterm{def:g9:inside-the-atom:proton}{protón} de \omterm{def:g11:functional-groups:carboxylic-acid}{éster}.
El triplete del etanol: \qty{3}{mm} para 3 \omterm{def:g9:inside-the-atom:proton}{protones}, \qty{1}{mm} por
\omterm{def:g9:inside-the-atom:proton}{protón} de etanol. Como cada \omterm{def:g7:atoms-and-molecules:molecule}{molécula} aporta un \ce{CH3} a estas señales,
las cantidades están en la razón $1/10$: alrededor de una \omterm{def:g7:atoms-and-molecules:molecule}{molécula} de
etanol por cada diez de \omterm{def:g11:functional-groups:carboxylic-acid}{éster}.
\end{solution}

\begin{solution}{exo:g12:proton-nmr:15}
Cuatro señales: los dos \ce{CH3} (6 \omterm{def:g9:inside-the-atom:proton}{protones}), un doblete; el \ce{CH} (1
\omterm{def:g9:inside-the-atom:proton}{protón}), desdoblado por 6 + 2 = 8 vecinos en 9 líneas; el \ce{CH2} (2
\omterm{def:g9:inside-the-atom:proton}{protones}), vecino del \ce{O}, un doblete, en el intervalo \ce{H-C-O}; el
\ce{OH} (1 \omterm{def:g9:inside-the-atom:proton}{protón}), un singlete. La señal del \ce{CH} es la que tiene más
líneas.
\end{solution}

\begin{solution}{pb:g12:proton-nmr:1}
\textbf{1.} Un \omterm{def:g11:organic-skeletons:alkane}{alcano} de cuatro carbonos es \ce{C4H10}; un \ce{C=O}
quita dos hidrógenos y los dos oxígenos no añaden ninguno: \ce{C4H8O2},
la fórmula de los ácidos y \omterm{def:g11:functional-groups:carboxylic-acid}{ésteres} de cuatro carbonos.

\textbf{2.} Ácido butanoico: \ce{CH3-CH2-CH2-COOH};\\
etanoato de etilo: \ce{CH3-CO-O-CH2-CH3};\\
propanoato de metilo: \ce{CH3-CH2-CO-O-CH3};\\
metanoato de propilo: \ce{H-CO-O-CH2-CH2-CH3}.

\textbf{3.} Los tres \omterm{def:g11:functional-groups:carboxylic-acid}{ésteres} comparten el grupo \omterm{def:g11:functional-groups:carboxylic-acid}{éster}; los cuatro tienen
un \ce{C=O}.

\textbf{4.} Un \ce{C=O}, en el intervalo de los \omterm{def:g11:functional-groups:carboxylic-acid}{ésteres} (cerca de \qty{1735}{cm^{-1}}).

\textbf{5.} La banda \ce{O-H} muy ancha, de 2500 a
\qty{3300}{cm^{-1}}: falta, así que no es el ácido butanoico.

\textbf{6.} No: los tres \omterm{def:g11:functional-groups:carboxylic-acid}{ésteres} tienen los mismos grupos.

\textbf{7.} Tres.

\textbf{8.} Un grupo etilo \ce{CH3-CH2-}: el \ce{CH2} desdoblado por 3
\omterm{def:g9:inside-the-atom:proton}{protones} (cuadruplete) y el \ce{CH3} por 2 (triplete).

\textbf{9.} Sus \omterm{def:g9:inside-the-atom:proton}{protones} no tienen \omterm{def:g9:inside-the-atom:proton}{protones} en los \omterm{def:g7:atoms-and-molecules:atom}{átomos} vecinos: un
\ce{CH3} unido al carbono del \ce{C=O} o al oxígeno del grupo \omterm{def:g11:functional-groups:carboxylic-acid}{éster}.

\textbf{10.} El \ce{CH2} está unido a un oxígeno (intervalo \ce{H-C-O}).

\textbf{11.} Un triplete (\ce{CH3}), un cuadruplete cerca de 2.0--\qty{2.4}{ppm}
(\ce{CH2} vecino del \ce{C=O}) y un singlete en el intervalo \ce{H-C-O}
(\ce{O-CH3}). El cuadruplete quedaría muy por debajo de \qty{4.12}{ppm}
y el singlete muy por encima de \qty{2.04}{ppm}: no es el compuesto
desconocido.

\textbf{12.} Cuatro señales (un singlete muy a la izquierda, un
triplete, una señal de 6 líneas y un triplete): no es el compuesto
desconocido, que tiene tres.

\textbf{13.} Etanoato de etilo, \ce{CH3-CO-O-CH2-CH3}: singlete a
\qty{2.04}{ppm} = \ce{CH3-C=O}; cuadruplete a \qty{4.12}{ppm} = \ce{O-CH2};
triplete a \qty{1.26}{ppm} = el \ce{CH3} del extremo.

\textbf{14.} Singlete 3, cuadruplete 2, triplete 3.

\textbf{15.} $72 / 8 = \qty{9}{mm}$ por \omterm{def:g9:inside-the-atom:proton}{protón}.

\textbf{16.} Singlete $3 \times 9 = \qty{27}{mm}$; triplete
\qty{27}{mm}.

\textbf{17.} $18 + 27 + 27 = \qty{72}{mm}$.

\textbf{18.} Sí: los \omterm{def:g11:functional-groups:carboxylic-acid}{ésteres} pequeños son conocidos por sus olores afrutados.

\textbf{19.} El escalón del cuadruplete mide \qty{18}{mm} de los \qty{72}{mm}.
\end{solution}
